\documentclass[11pt]{article}

\usepackage[T1]{fontenc}
\usepackage[utf8]{inputenc}
\usepackage{lmodern}
\usepackage[margin=1.1in]{geometry}
\usepackage{amsmath,amssymb,amsthm,mathtools}
\usepackage{enumitem}
\usepackage{microtype}
\usepackage{xcolor}
\usepackage{array}
\usepackage{titlesec}
\usepackage{hyperref}

\titleformat{\section}
  {\normalfont\LARGE\bfseries}{\thesection}{1em}{}
\titleformat{\subsection}
  {\normalfont\Large\bfseries}{\thesubsection}{1em}{}

\hypersetup{
  pdftitle={Generic Degree Three in JC(3): Collision Geometry and Galois Group S3},
  pdfauthor={Chloe van der Vlugt},
  colorlinks=true,
  linkcolor=blue!55!black,
  citecolor=green!40!black,
  urlcolor=blue!60!black
}

\newtheorem{theorem}{Theorem}[section]
\newtheorem{proposition}[theorem]{Proposition}
\newtheorem{corollary}[theorem]{Corollary}
\theoremstyle{remark}
\newtheorem*{remark}{Remark}

\DeclareMathOperator{\Spec}{Spec}
\DeclareMathOperator{\Frac}{Frac}
\DeclareMathOperator{\Gal}{Gal}
\DeclareMathOperator{\Hom}{Hom}
\DeclareMathOperator{\Norm}{Norm}
\DeclareMathOperator{\core}{core}

\newcommand{\A}{\mathbb A}
\newcommand{\C}{\mathbb C}
\newcommand{\Obs}{\operatorname{Obs}}
\newcommand{\id}{\mathrm{id}}
\newcommand{\leanpath}[1]{\texttt{\nolinkurl{#1}}}

\newenvironment{leanfiles}
  {\par\medskip\begingroup\footnotesize\raggedright\noindent
   \textit{Lean correspondence.}\ }
  {\par\endgroup\medskip}

\title{\texorpdfstring
  {Generic Degree Three in \(JC(3)\): Collision Geometry\\
   and Galois Group \(S_3\)}
  {Generic Degree Three in JC(3): Collision Geometry and Galois Group S3}}
\author{Chloe van der Vlugt}
\date{\today}

\begin{document}

\maketitle

\begin{abstract}
We study possible counterexamples to the complex three-dimensional
Jacobian conjecture.  For a polynomial map
\(F\colon\A^n_{\C}\to\A^n_{\C}\), let \(R_F\) be the scheme of pairs
\((x,y)\) satisfying \(F(x)=F(y)\).  In
\(S=\C[x_1,\dots,x_n,y_1,\dots,y_n]\), the collision scheme and its
diagonal are defined respectively by
\[
  I_R(F)=\bigl(F_i(x)-F_i(y)\bigr)_{i=1}^n,
  \qquad
  I_\Delta=(x_i-y_i)_{i=1}^n.
\]
Since \(I_R(F)\subseteq I_\Delta\), the quotient
\[
  \Obs(F):=I_\Delta/I_R(F)\subseteq C_F:=S/I_R(F)
\]
defines the diagonal inside \(R_F\).  Over \(\C\), the map \(F\) is a
polynomial automorphism precisely when \(I_R(F)=I_\Delta\), equivalently
when \(\Obs(F)=0\).  The coordinate functions of \(F\) generate
\[
  B=\C[F_1,\dots,F_n]\subseteq A=\C[x_1,\dots,x_n],
\]
and hence a function-field extension
\(K=\Frac(B)\subseteq L=\Frac(A)\).
When \(n=3\) and \([L:K]=3\), the map \(F\) is a counterexample to
\(JC(3)\) precisely when it is Keller.

For every such counterexample, the generic collision algebra decomposes as
\[
  K\otimes_B C_F\simeq_K L\times N,
\]
where \(N\) is the normal closure of \(L/K\).  Consequently, the normal
closure of the function-field extension attached to every
generic-degree-three counterexample to \(JC(3)\) has Galois group \(S_3\).
\end{abstract}

\begin{center}
  \small Licensed under
  \href{https://creativecommons.org/licenses/by/4.0/}
    {CC BY 4.0}.
\end{center}

\tableofcontents

\section{Introduction}
\label{sec:introduction}

The complex Jacobian conjecture \(JC(n)\) asks whether every polynomial map
\[
  F=(F_1,\dots,F_n)\colon\A^n_{\C}\longrightarrow\A^n_{\C},
  \qquad \det JF\in\C^\times,
\]
has a polynomial inverse.  Such an \(F\) is called a Keller map.  The
Jacobian condition makes \(F\) an \'etale morphism; the conjecture asks
whether this local condition forces global invertibility.

Put
\[
  S=\C[x_1,\dots,x_n,y_1,\dots,y_n].
\]
The equations
\[
  F_i(x)-F_i(y)=0\qquad(1\leq i\leq n)
\]
generate the collision ideal, while the equations of the diagonal generate
the diagonal ideal:
\[
  I_R(F)=\bigl(F_i(x)-F_i(y)\bigr)_{i=1}^n,
  \qquad
  I_\Delta=(x_i-y_i)_{i=1}^n.
\]
Thus the collision scheme is
\[
  R_F=\Spec C_F,
  \qquad
  C_F=S/I_R(F).
\]
Its complex points are precisely the pairs \((x,y)\) with \(F(x)=F(y)\);
equivalently, \(R_F\) is the fiber square of \(F\).  It contains the
diagonal \(\Delta\), consisting of the pairs \((x,x)\).  Since every
collision equation vanishes on the diagonal,
\[
  I_R(F)\subseteq I_\Delta.
\]
Its image in \(C_F\) is the ideal
\[
  \Obs(F):=I_\Delta/I_R(F)\subseteq C_F,
\]
which defines the diagonal inside \(R_F\).  By Ax--Grothendieck,
\[
  I_R(F)=I_\Delta
  \quad\Longleftrightarrow\quad
  \Obs(F)=0
  \quad\Longleftrightarrow\quad
  R_F=\Delta
  \quad\Longleftrightarrow\quad
  F\text{ is a polynomial automorphism}.
\]
Thus, among Keller maps, strict containment
\(I_R(F)\subsetneq I_\Delta\), equivalently \(\Obs(F)\neq0\), is precisely
the counterexample condition.

We organize this possible strict containment by generic degree.  On
coordinate rings, \(F\) induces
\[
  F^*\colon\C[t_1,\dots,t_n]\longrightarrow A,
  \qquad t_i\longmapsto F_i.
\]
Let
\[
  B=\operatorname{im}(F^*)=\C[F_1,\dots,F_n]\subseteq A,
  \qquad K=\Frac(B),
  \qquad L=\Frac(A).
\]
For a Keller map, the \(F_i\) are algebraically independent, so \(F^*\)
identifies the target coordinate ring with \(B\).  Thus
\(K=\C(F_1,\dots,F_n)\subseteq L=\C(x_1,\dots,x_n)\) is the induced
finite separable extension of function fields, and \([L:K]\) is the
generic degree.  If \(L/K\) is normal, the classical Galois case of the
Jacobian conjecture implies that \(F\) is a polynomial automorphism.  Since
extensions of degree one and separable extensions of degree two are normal,
neither degree can occur for a counterexample.  Generic degree three is
therefore the first possible counterexample degree.

We henceforth study polynomial maps
\(F\colon\A^3_{\C}\to\A^3_{\C}\) with generic degree \([L:K]=3\).
Here \(JC(3)\) denotes the conjecture in three variables, while
\([L:K]=3\) is the separate generic-degree condition.  Since polynomial
automorphisms have generic degree one, every map with \([L:K]=3\) is already
nonautomorphic.  In this generic-degree-three case,
\[
  F\text{ is a counterexample to }JC(3)
  \quad\Longleftrightarrow\quad
  F\text{ is Keller}.
\]
Once \([L:K]=3\) is assumed, the residual decomposition is field-theoretic.
The generic collision algebra has a diagonal factor \(L\) and a quadratic
residual factor \(D_\alpha\):
\[
  K\otimes_BC_F\simeq_K L\times D_\alpha.
\]
The quadratic residual factor is split in the normal branch and is a field
in the nonnormal branch.  In the latter case it is \(L\)-isomorphic to the
normal closure \(N\).  The Keller condition identifies the counterexample
locus and, through Galois rigidity, excludes the normal branch.  Hence every
generic-degree-three counterexample to \(JC(3)\) has
\[
  K\otimes_BC_F\simeq_K L\times N,
  \qquad
  \Gal(N/K)\cong S_3.
\]
Section~\ref{sec:general-construction} develops the dimension-independent
collision obstruction, generic base change, marked-root decomposition, and
automorphism criterion.  Section~\ref{sec:cubic-s3-obstruction} specializes
these constructions to cubic extensions.  Section
\ref{sec:further-directions} isolates the normalization and boundary
questions leading to the planar sequel.  The closing low-degree corollary
records the precise handoff to that paper.

\section{General construction}
\label{sec:general-construction}

For a fixed dimension \(n\), where can \(JC(n)\) fail?  We begin with an
arbitrary polynomial map.  Its collision ideal is canonically contained in
the diagonal ideal, and \(\Obs(F)=I_\Delta/I_R(F)\) measures their failure to
agree.  Once \(F\) is assumed Keller, strict containment is exactly the
counterexample condition.  No Keller hypothesis enters the collision
construction.

\subsection{The two ideals and their obstruction}
\label{sec:collision-ideals}

Throughout the paper, the base field is \(\C\).  Let
\[
  F=(F_1,\dots,F_n)\colon
  X=\A^n_{\C}\longrightarrow\A^n_{\C}
\]
be a polynomial map, and let
\[
  A=\C[x_1,\dots,x_n],
  \qquad
  S=\C[x_1,\dots,x_n,y_1,\dots,y_n].
\]
Define
\[
  I_R(F)=\bigl(F_i(x)-F_i(y)\bigr)_{i=1}^{n},
  \qquad
  I_\Delta=(x_i-y_i)_{i=1}^{n},
  \qquad
  C_F:=S/I_R(F),
  \qquad
  R_F:=\Spec C_F.
\]
For a fixed map \(F\), we abbreviate \(I_R(F)\) to \(I_R\).

\begin{proposition}[Canonical containment]
\label{prop:collision-contained-in-diagonal}
For every polynomial map \(F\),
\[
  I_R(F)\subseteq I_\Delta.
\]
\end{proposition}

\begin{proof}
For \(P\in\C[x_1,\dots,x_n]\), the difference \(P(x)-P(y)\) belongs to
\((x_1-y_1,\dots,x_n-y_n)\).  This follows by induction on \(P\), or by
telescoping each monomial one coordinate at a time.  Apply the statement
to each coordinate \(F_i\).
\end{proof}

Let \(q_F\colon S\twoheadrightarrow C_F\) be the quotient map.  Diagonal
substitution \(S\twoheadrightarrow A\), given by
\(s(x,y)\mapsto s(x,x)\), has kernel \(I_\Delta\).  Since
\(I_R\subseteq I_\Delta\), it descends uniquely to a surjection
\[
  \bar\mu_F\colon C_F\twoheadrightarrow A,
  \qquad
  [s]_{I_R}\longmapsto s(x,x).
\]
We define the ideal of the diagonal in \(R_F\) by
\[
  \Obs(F):=q_F(I_\Delta)=I_\Delta/I_R\subseteq C_F.
\]

\begin{proposition}[The diagonal kernel]
\label{prop:obstruction-kernel}
As ideals of \(C_F\),
\[
  \ker(\bar\mu_F)
  =
  \Obs(F)
  =
  I_\Delta/I_R.
\]
Equivalently,
\[
  \ker(\bar\mu_F)=\Obs(F)=0
  \quad\Longleftrightarrow\quad
  I_R=I_\Delta.
\]
\end{proposition}

\begin{proof}
The construction of \(\bar\mu_F\) identifies its kernel with
\(I_\Delta/I_R=\Obs(F)\).  Since \(I_R\subseteq I_\Delta\),
\[
  I_\Delta/I_R=0
  \quad\Longleftrightarrow\quad
  I_R=I_\Delta.
\]
\end{proof}

\subsection{The fiber-product interpretation}
\label{sec:fiber-product}

Let
\[
  A=\C[x_1,\dots,x_n],\qquad
  B=\C[F_1,\dots,F_n]\subseteq A.
\]
Let \(X=\Spec A\), let \(Y=\Spec B\), and let
\[
  f_F\colon X\longrightarrow Y
\]
be the morphism induced by \(B\hookrightarrow A\).
The homomorphisms
\(\iota_x,\iota_y\colon A\rightrightarrows C_F\), defined by
\(\iota_x(a)=[a(x)]\) and \(\iota_y(a)=[a(y)]\), agree on \(B\) because
\([F_i(x)]=[F_i(y)]\) in \(C_F\).  The equality
\(\iota_x|_B=\iota_y|_B\) defines the \(B\)-algebra structure on \(C_F\).

\begin{theorem}[Fiber-product interpretation]
\label{thm:collision-tensor-product}
There is a canonical isomorphism
\[
  R_F\cong X\times_YX
\]
under which \(\Spec(\bar\mu_F)\colon X\to R_F\) is the relative diagonal
\(\Delta_{X/Y}\).  Dually, there is a canonical \(B\)-algebra isomorphism
\[
  C_F\xrightarrow{\sim}_B A\otimes_BA
\]
under which \(\bar\mu_F\) is tensor multiplication
\(a\otimes a'\mapsto aa'\).
\end{theorem}

\begin{proof}
For every commutative \(\C\)-algebra \(T\), maps from either \(C_F\) or
\(A\otimes_BA\) to \(T\) are pairs
\(f,g\colon A\rightrightarrows T\) agreeing on \(B\).  In the presentation
of \(C_F\), this agreement is imposed by the relations
\(F_i(x)=F_i(y)\), since the \(F_i\) generate \(B\).  The common universal
property gives the algebra isomorphism, and applying \(\Spec\) gives the
scheme isomorphism.  Tensor multiplication corresponds to
\(\bar\mu_F\), so its dual is the relative diagonal.
\end{proof}

On complex points,
\[
  R_F(\C)=\{(u,v):F(u)=F(v)\}.
\]
Inside \(X\times X=\Spec S\), the ideals \(I_R\) and \(I_\Delta\)
define \(X\times_YX\) and its diagonal.  Thus
\[
  I_R=I_\Delta
  \quad\Longleftrightarrow\quad
  X\times_YX=\Delta_{X/Y}(X)
\]
as closed subschemes.  This is the geometric form of obstruction
vanishing.

\subsection{Generic fibers and marked-root decomposition}
\label{sec:generic-galois-bridge}

Let \(F\) be a polynomial map.  Retain \(A\), and set
\[
  B=\C[F_1,\dots,F_n]\subseteq A.
\]
Let
\[
  K=\Frac(B),
  \qquad
  L=\Frac(A).
\]

\begin{proposition}[Keller generic geometry]
\label{prop:keller-generic-geometry}
If \(F\) is Keller, then its coordinate functions are algebraically
independent, the morphism
\[
  f_F\colon\Spec A\longrightarrow\Spec B
\]
is \'etale and quasi-finite, and \(L/K\) is a finite separable extension.
Consequently, every Keller map, and in particular every complex
Jacobian-conjecture counterexample, has a finite separable generic field
extension.
\end{proposition}

\begin{proof}
The identity
\[
  dF_1\wedge\cdots\wedge dF_n
  =\det(JF)\,dx_1\wedge\cdots\wedge dx_n\neq0
\]
gives algebraic independence, so
\(B\cong\C[u_1,\dots,u_n]\).  As a \(B\)-algebra, \(A\) has the
presentation
\[
  A\cong
  \C[u_1,\dots,u_n,X_1,\dots,X_n]
  \big/\bigl(F_i(X)-u_i\bigr)_{i=1}^n.
\]
Its relative Jacobian determinant is the unit \(\det(JF)\).  The Jacobian
criterion therefore makes \(f_F\) \'etale, hence quasi-finite.  Since
\(B\subseteq A\), the morphism is dominant, so quasi-finiteness gives
\([L:K]<\infty\).  Its generic point is \'etale, and therefore \(L/K\) is
separable.
\end{proof}

For the remainder of this subsection, assume that \(F\) is dominant and
generically finite.  Let
\[
  \theta\colon K\otimes_BA\longrightarrow L,
  \qquad
  \lambda\otimes a\longmapsto\lambda a.
\]
Let
\[
  \mu_L\colon L\otimes_KL\longrightarrow L,
  \qquad
  \ell\otimes\ell'\longmapsto\ell\ell',
\]
denote tensor multiplication.

\begin{theorem}[Generic collision bridge]
\label{thm:generic-collision-base-change}
The map \(\theta\) is an isomorphism.  Consequently, there is a
\(K\)-algebra isomorphism
\[
  \Phi_F\colon K\otimes_BC_F
  \xrightarrow{\sim}_K
  L\otimes_KL,
\]
and the following square commutes:
\[
  \begin{array}{ccc}
    K\otimes_BC_F
      &\xrightarrow{\ 1\otimes\bar\mu_F\ }&
    K\otimes_BA\\
    {\scriptstyle\Phi_F}\downarrow
      &&\downarrow{\scriptstyle\theta\ \sim}\\
    L\otimes_KL
      &\xrightarrow{\ \mu_L\ }&
    L.
  \end{array}
\]
\end{theorem}

\begin{proof}
Let \(U=B\setminus\{0\}\).  Since \(K=U^{-1}B\), localization gives
\[
  K\otimes_BA\cong U^{-1}A\subseteq L,
\]
and under this identification \(\theta\) is the inclusion.  Since \(L/K\)
is finite, each \(x_i\) is algebraic over \(K\); hence
\[
  U^{-1}A=K[x_1,\dots,x_n]
\]
is a finite-dimensional domain over \(K\), and therefore a field.  It
contains \(A\), so it contains \(\Frac(A)=L\).  Thus \(U^{-1}A=L\), and
\(\theta\) is an isomorphism.

Using Theorem~\ref{thm:collision-tensor-product} and standard tensor base
change, define \(\Phi_F\) by the composite
\[
  K\otimes_BC_F
  \cong
  K\otimes_B(A\otimes_BA)
  \cong
  (K\otimes_BA)\otimes_K(K\otimes_BA)
  \xrightarrow{\ \theta\otimes\theta\ }
  L\otimes_KL.
\]
Under this composite, both routes around the square send
\(\lambda\otimes(a\otimes a')\) to \(\lambda aa'\).  Hence the square
commutes.
\end{proof}

Suppose now that \(L/K\) is separable and has a power basis
\(L=K(\alpha)\).  Let \(m_\alpha(T)\in K[T]\) be the minimal polynomial
of \(\alpha\), and let \(h_\alpha(T)\in L[T]\) be determined by
\[
  m_\alpha(T)=(T-\alpha)h_\alpha(T)
  \qquad\text{in }L[T].
\]
Set
\[
  D_\alpha:=L[T]/(h_\alpha).
\]

\begin{theorem}[Marked-root tensor decomposition]
\label{thm:marked-root-tensor-decomposition}
Give \(L\otimes_KL\) its \(L\)-algebra structure through the left tensor
factor.  There is an \(L\)-algebra isomorphism
\[
  \Phi_\alpha\colon L\otimes_KL
  \xrightarrow{\sim}_L
  L\times D_\alpha,
\]
under which \(\mu_L\) is the first-coordinate map.
\end{theorem}

\begin{proof}
View \(m_\alpha\) in \(L[T]\).  Separability gives
\(h_\alpha(\alpha)=m_\alpha'(\alpha)\neq0\), so \(T-\alpha\) and
\(h_\alpha\) are coprime.  The power-basis isomorphism and the Chinese
remainder theorem give the composite
\[
  L\otimes_KL
  \xrightarrow{\sim}_L L[T]/(m_\alpha)
  \xrightarrow{\sim}_L
  L[T]/(T-\alpha)\times L[T]/(h_\alpha)
  \xrightarrow{\sim}_L L\times D_\alpha.
\]
The first coordinate of the composite is evaluation at \(T=\alpha\),
which is precisely \(\mu_L\).
\end{proof}

Together with Theorem~\ref{thm:generic-collision-base-change}, this gives
\[
  K\otimes_BC_F
  \simeq_K
  L\times D_\alpha,
\]
with the base-changed diagonal evaluation equal to the first-coordinate
map.

\subsection{Automorphism detection}
\label{sec:automorphism-detection}

\begin{theorem}[Automorphism criterion]
\label{thm:automorphism-criterion}
For an arbitrary polynomial map
\(F\colon\A^n_{\C}\to\A^n_{\C}\), the following are equivalent:
\begin{enumerate}[label=\textup{(\arabic*)}]
  \item \(F\) is a polynomial automorphism;
  \item \(F\) is injective on \(\A^n(\C)\);
  \item \(I_R(F)=I_\Delta\);
  \item \(\Obs(F)=0\).
\end{enumerate}
\end{theorem}

\begin{proof}
Suppose first that \(F\) has polynomial inverse
\(G=(G_1,\dots,G_n)\).  Then
\[
  x_i-y_i
  =
  G_i(F(x))-G_i(F(y))
  \in I_R
\]
for every \(i\), so \(I_\Delta\subseteq I_R\).  Proposition
\ref{prop:collision-contained-in-diagonal} gives equality.  Thus
\textup{(1)} implies \textup{(3)}.

Suppose \(I_R=I_\Delta\), and let
\(\operatorname{ev}_{(a,b)}\colon S\to\C\) be evaluation at
\((a,b)\).  If \(F(a)=F(b)\), then
\[
  I_R\subseteq\ker(\operatorname{ev}_{(a,b)}).
\]
Therefore \(x_i-y_i\in I_\Delta=I_R\) gives \(a_i-b_i=0\) for every
\(i\).  Thus \(a=b\), proving
\textup{(3)}\(\Rightarrow\)\textup{(2)}.  The
Ax--Grothendieck theorem \cite{Ax1969} gives
\textup{(2)}\(\Rightarrow\)\textup{(1)}.  Finally,
Proposition~\ref{prop:obstruction-kernel} gives
\textup{(3)}\(\Longleftrightarrow\)\textup{(4)}.
\end{proof}

Applied to Keller maps, Theorem~\ref{thm:automorphism-criterion} says
that \(JC(n)\) is equivalent to universal vanishing of \(\Obs(F)\), or
equivalently to \(I_R(F)=I_\Delta\).

We will also use the classical Galois case of the conjecture.  It is the
step that rules out the normal cubic branch under the Keller hypothesis.

\begin{theorem}[Classical Galois rigidity]
\label{thm:classical-galois-rigidity}
Let
\[
  F\colon\A^n_{\C}\longrightarrow\A^n_{\C},
  \qquad \det JF\in\C^\times,
\]
and put
\[
  K=\C(F_1,\dots,F_n),
  \qquad
  L=\C(x_1,\dots,x_n).
\]
If \(L/K\) is normal, then \(F\) is a polynomial automorphism and
\(L=K\).
\end{theorem}

\begin{proof}
The Keller condition makes \(L/K\) finite and separable.  Thus normality
makes it Galois, and the conclusion is the classical Galois case of the
Jacobian conjecture due to Campbell, with algebraic treatments by Razar
and Wright \cite{Campbell1973,Razar1979,Wright1981}.
\end{proof}

\begin{leanfiles}
\leanpath{CollisionIdeals/General/Collision/Diagonal.lean},
\leanpath{CollisionIdeals/General/Collision/Ideal.lean},
\leanpath{CollisionIdeals/General/Collision/CollisionDiagonal.lean},
\leanpath{CollisionIdeals/General/FiberProduct/Polynomial.lean},
\leanpath{CollisionIdeals/General/GenericFiber/FunctionField.lean},
\leanpath{CollisionIdeals/General/GenericFiber/CollisionDiagonal.lean},
\leanpath{CollisionIdeals/General/GenericFiber/MarkedRootDecomposition.lean},
\leanpath{CollisionIdeals/General/Keller/Interfaces.lean},
\leanpath{CollisionIdeals/General/Automorphism/Equivalences.lean}, and
\leanpath{CollisionIdeals/General/Automorphism/GaloisRigidity.lean}.
The polynomial fiber-product module uses the abstract target polynomial
ring; the equivalence over the actual image algebra
\(B=\C[F_1,\ldots,F_n]\) and its generic base change are owned by the
generic-function-field module.  The Lean form of
Theorem~\ref{thm:generic-collision-base-change} takes surjectivity of
\(\theta\) as a parameter; the proof above derives it from generic
finiteness.  The predicates \leanpath{KellerEtaleBridge} and
\leanpath{KellerFlatBridge} expose parts of Proposition
\ref{prop:keller-generic-geometry} as theorem targets; Lean does not yet
derive them, generic finiteness, or separability from \leanpath{IsKeller}.
Ax--Grothendieck also remains an explicit parameter of the Lean
automorphism equivalence.
\end{leanfiles}

% The planar specialization is developed in the companion Paper II.

\section{\texorpdfstring{The cubic \(S_3\) obstruction in dimension three}
  {The cubic S3 obstruction in dimension three}}
\label{sec:cubic-s3-obstruction}

We now specialize the dimension-independent constructions of
Section~\ref{sec:general-construction} to generic degree three.  After the
diagonal is exhibited as the first factor, the complementary
generic collision algebra has degree two over \(L\).  For a map of generic
degree three it is either split,
corresponding to a normal cubic extension, or connected, corresponding to
a nonnormal cubic extension.  Degree three already rules out polynomial
automorphy, so the counterexamples to \(JC(3)\) in generic degree three are
precisely the Keller maps.  For a Keller map, however, the split branch is
impossible: it makes \(L/K\) normal and hence, by
Theorem~\ref{thm:classical-galois-rigidity}, forces \(F\) to be a polynomial
automorphism, contradicting \([L:K]=3\).  Hence only the connected branch
remains; its residual factor is the normal closure, with Galois group
\(S_3\).

\subsection{The cubic extension and its normal closure}
\label{sec:cubic-extension}

Let
\[
  F=(F_1,F_2,F_3)\colon\A^3_{\C}\longrightarrow\A^3_{\C}
\]
be dominant and generically finite.  Let
\[
  A=\C[x_1,x_2,x_3],
  \qquad
  B=\C[F_1,F_2,F_3]\subseteq A,
  \qquad
  K=\Frac(B),
  \qquad
  L=\Frac(A).
\]
Assume that \(F\) has generic degree three:
\[
  [L:K]=3.
\]
Because \(K\) has characteristic zero, \(L/K\) is separable.  Choose a
primitive element \(\alpha\)
\cite{StacksPrimitiveElement},
so that
\(L=K(\alpha)\) with power basis
\(1,\alpha,\alpha^2\).
Retain the notation of
Theorem~\ref{thm:marked-root-tensor-decomposition}: if \(m_\alpha\in K[T]\)
is the minimal polynomial of \(\alpha\), then
\[
  m_\alpha(T)=(T-\alpha)h_\alpha(T)
  \qquad\text{in }L[T].
\]
The residual factor in the marked-root decomposition is
\[
  D_\alpha=L[T]/(h_\alpha).
\]
Since \([L:K]=3\), it has dimension two over \(L\).
Let \(N/K\) be the normal closure of \(L/K\), containing \(L\) through
the marked inclusion.  Since \(L/K\) is finite separable, \(N/K\) is
Galois
\cite{StacksNormalClosureGalois}.
Let
\[
  G=\Gal(N/K),
  \qquad
  H=\Gal(N/L)=\{g\in G:g|_L=\id_L\},
\]
where \(H\) fixes the marked copy of \(L\) in \(N\).
Restriction of automorphisms gives a bijection
\[
  G/H\xrightarrow{\sim}\Hom_K(L,N),
  \qquad gH\longmapsto g|_L.
\]
Because the conjugates of \(L\) generate its normal closure \(N\), an
element of \(G\) fixing every conjugate embedding fixes \(N\).  Hence the
coset action is faithful and
\[
  \core_G(H)=\bigcap_{g\in G}gHg^{-1}=1.
\]

The residual collision algebra detects which cubic branch occurs.  We
include the short argument in order to fix its marked identification with
the normal closure.

\begin{proposition}[Connected cubic residual factor]
\label{prop:cubic-tensor-decomposition}
The following conditions are equivalent:
\begin{enumerate}[label=\textup{(\roman*)}]
  \item \(\Spec D_\alpha\) is connected;
  \item \(D_\alpha\) is a field;
  \item \(h_\alpha\) is irreducible over \(L\);
  \item \(L/K\) is nonnormal.
\end{enumerate}
When they hold, there is an \(L\)-algebra isomorphism
\[
  \varphi_\alpha\colon D_\alpha\xrightarrow{\sim}_L N.
\]
Consequently,
\[
  \Psi_\alpha\colon L\otimes_KL\xrightarrow{\sim}_L L\times N,
\]
under which \(\mu_L\) is the first-coordinate map.
\end{proposition}

\begin{proof}
Write
\[
  m_\alpha(T)=T^3+a_2T^2+a_1T+a_0.
\]
The quadratic \(h_\alpha\) is separable, so \(D_\alpha\) is finite
\'etale over \(L\); finite \'etale algebras over a field are finite
products of finite separable field extensions
\cite{StacksEtaleAlgebrasField}.
If \(h_\alpha\) is irreducible, then
\(D_\alpha\) is a field and hence has connected spectrum.  If it is
reducible, it is a product of two distinct linear factors; the Chinese
remainder theorem then gives \(D_\alpha\simeq_L L\times L\), whose spectrum
is disconnected.  This proves the equivalence of \textup{(i)},
\textup{(ii)}, and \textup{(iii)}.

If \(h_\alpha\) is reducible, it has a root \(\beta\in L\).
Separability gives \(\beta\neq\alpha\), and the third root is
\[
  \gamma=-a_2-\alpha-\beta\in L.
\]
Thus \(m_\alpha\) would split over \(L\).  Since
\(L=K(\alpha)\), the standard generator criterion for normality
\cite{StacksNormalGeneratorCriterion}
shows that \(L/K\) would be normal.  Conversely,
if \(L/K\) is normal, then \(m_\alpha\) splits over \(L\), so
\(h_\alpha\) is reducible.  Hence \textup{(iii)} and \textup{(iv)} are
equivalent.

Assume these equivalent conditions, and let \(\beta\) denote the image of
\(T\) in \(D_\alpha\).  This field contains the three roots
\[
  \alpha,\qquad \beta,\qquad -a_2-\alpha-\beta
\]
of \(m_\alpha\), and is generated over \(L\) by \(\beta\).  It is
therefore the splitting field of \(m_\alpha\), which is the normal
closure \(N\) of \(L/K\)
\cite{StacksNormalClosureExistence}.
Choose the resulting \(L\)-algebra
isomorphism
\[
  \varphi_\alpha\colon D_\alpha\xrightarrow{\sim}_L N.
\]

Let
\[
  \Phi_\alpha\colon
  L\otimes_KL
  \xrightarrow{\sim}_L
  L\times D_\alpha
\]
be the marked-root decomposition of
Theorem~\ref{thm:marked-root-tensor-decomposition}.  Under this
isomorphism, \(\mu_L\) is the first-coordinate map.
Composing its second factor with \(\varphi_\alpha\), let
\[
  \Psi_\alpha=(\id_L\times\varphi_\alpha)\circ\Phi_\alpha.
\]
Since \(\id_L\times\varphi_\alpha\) leaves the first factor unchanged,
the same description of \(\mu_L\) holds under \(\Psi_\alpha\).
\end{proof}

\begin{corollary}[The generic-degree-three counterexample locus]
\label{cor:cubic-counterexample-locus}
The map \(F\) is not a polynomial automorphism.  Equivalently, on this
generic-degree-three locus,
\[
  F\text{ is a counterexample to }JC(3)
  \quad\Longleftrightarrow\quad
  F\text{ is Keller}.
\]
When these equivalent conditions hold, \(L/K\) is nonnormal and
\(\Spec D_\alpha\) is connected.
\end{corollary}

\begin{proof}
If \(F\) were a polynomial automorphism, its coordinate functions would
generate \(A\), so \(B=A\) and \(L=K\), contrary to \([L:K]=3\).
The displayed equivalence now follows from the definition of a
counterexample to \(JC(3)\).

Assume these conditions, so that \(F\) is Keller.  If \(L/K\) were
normal, Theorem~\ref{thm:classical-galois-rigidity} would give \(L=K\),
again contradicting \([L:K]=3\).  Proposition
\ref{prop:cubic-tensor-decomposition} now identifies nonnormality with
connectedness of \(\Spec D_\alpha\).
\end{proof}

\begin{theorem}[Generic cubic \(S_3\) collision]
\label{thm:cubic-s3-collision}
Assume that \(\Spec D_\alpha\) is connected.  There is a \(K\)-algebra
isomorphism
\[
  \Psi_F\colon K\otimes_BC_F\xrightarrow{\sim}_K L\times N
\]
under which, after the identification \(K\otimes_BA\simeq_K L\) of
Theorem~\ref{thm:generic-collision-base-change}, the base-changed diagonal
evaluation
\[
  1\otimes\bar\mu_F
  \colon K\otimes_BC_F\longrightarrow K\otimes_BA
\]
is the first-coordinate map.  In particular,
\[
  \Psi_F\!\left(
    \ker(1\otimes\bar\mu_F)
  \right)
  =0\times N,
  \qquad
  \Gal(N/K)\cong S_3,
\]
and
\[
  \Obs(F)\neq0,
  \qquad
  I_R\subsetneq I_\Delta.
\]
By Corollary~\ref{cor:cubic-counterexample-locus}, the connectedness
hypothesis holds whenever \(F\) is Keller; in generic degree three these
Keller maps are precisely the counterexamples to \(JC(3)\).
\end{theorem}

\begin{proof}
Proposition~\ref{prop:cubic-tensor-decomposition} gives an \(L\)-algebra
isomorphism
\[
  \Psi_\alpha\colon L\otimes_KL\xrightarrow{\sim}_L L\times N
\]
under which \(\mu_L\) is the first-coordinate map.  Theorem
\ref{thm:generic-collision-base-change} gives a \(K\)-algebra isomorphism
\[
  \Phi_F\colon K\otimes_BC_F\xrightarrow{\sim}_K L\otimes_KL
\]
which carries the base-changed diagonal evaluation to \(\mu_L\).
After restricting \(\Psi_\alpha\) from \(L\) to \(K\), let
\[
  \Psi_F=\Psi_\alpha\circ\Phi_F.
\]
The asserted description of diagonal evaluation follows, and hence
\[
  \Psi_F\!\left(
    \ker(1\otimes\bar\mu_F)
  \right)
  =0\times N\neq0.
\]
If \(\bar\mu_F\colon C_F\to A\) were injective, flat base change along
the localization \(B\to K\) would make \(1\otimes\bar\mu_F\) injective.
This would contradict the displayed kernel.  Thus
\(\ker(\bar\mu_F)\neq0\), so Proposition
\ref{prop:obstruction-kernel} gives \(\Obs(F)\neq0\), and Proposition
\ref{prop:collision-contained-in-diagonal} gives
\(I_R\subsetneq I_\Delta\).

For the Galois group, Proposition
\ref{prop:cubic-tensor-decomposition} identifies \(N\) with the
degree-two field \(D_\alpha\) over \(L\).  The fixed-field degree formula
therefore gives
\[
  |H|=[N:L]=2,
  \qquad
  [G:H]=[L:K]=3.
\]
The normal-closure identities recalled above give \(\core_G(H)=1\).
Thus the action on the three cosets of \(H\) embeds \(G\) into \(S_3\).
But
\[
  |G|=[G:H]|H|=3\cdot2=6,
\]
so this embedding is an isomorphism \(G\cong S_3\).  We use here the
fixed-field degree formula from the fundamental theorem of Galois theory
\cite{StacksFundamentalGaloisTheory}.
\end{proof}

\begin{remark}[Ordered roots]
Let
\[
  m_\alpha(T)=T^3+a_2T^2+a_1T+a_0
\]
have roots \(\alpha_1,\alpha_2,\alpha_3\) in \(N\).  For distinct
\(i,j,k\),
\[
  \alpha_k=-a_2-\alpha_i-\alpha_j,
  \qquad
  K(\alpha_i,\alpha_j)
  =
  K(\alpha_1,\alpha_2,\alpha_3)
  =
  N.
\]
Thus the residual factor, which marks two distinct roots, is already
the ordered-root normal closure.
\end{remark}

\begin{leanfiles}
The cubic formalization is collected in
\leanpath{CollisionIdeals/ComplexThree/Cubic.lean}, using the generic
collision descent in
\leanpath{CollisionIdeals/General/GenericFiber/DiagonalObstruction.lean}
and the normal-closure API in
\leanpath{CollisionIdeals/General/Galois/NormalClosure.lean}.
From supplied separability, generic-source surjectivity, a cubic power
basis, marked normal-closure data, and nonnormality, Lean constructs
\(D_\alpha\simeq_L N\), proves \(H\neq1\) and
\(\Gal(N/K)\cong S_3\), and descends the residual factor to
\(\Obs(F)\neq0\).  The Keller hypothesis is then used only to package the
counterexample statement.  General inputs from the literature are isolated
as theorem interfaces rather than reproved in Lean; in particular,
classical Keller--Galois rigidity is recorded as
\leanpath{ComplexKellerGaloisRigidity}.  Detailed object-level parity is
maintained in \leanpath{SEMANTIC-PARITY.md}.
\end{leanfiles}

\section{Further directions}
\label{sec:further-directions}

The normal closure also carries a dimension-independent boundary picture,
although none of it is needed for the cubic collision theorem.  Let
\(L/K\) be finite separable, let \(N/K\) be its normal closure, and let
\(Y\) be a normal integral model with function field \(K\).  Write
\[
  G=\Gal(N/K),\qquad H=\Gal(N/L),\qquad
  Z=\Norm_N(Y)\longrightarrow \overline X=\Norm_L(Y)\longrightarrow Y.
\]
Suppose that a normal model \(X\) of \(L\) is an open subscheme of
\(\overline X\) over \(Y\), and that \(X\to Y\) is \'etale.

\begin{proposition}[Hidden inertia on a conjugate sheet]
\label{prop:generic-hidden-inertia}
Let \(E\in Z^{(1)}\) be ramified over \(Y\), with divisorial valuation
\(w_E\), restriction \(v=w_E|_K\), and inertia group \(I_E\).  Then some
\(g\in G\) satisfies
\[
  e\bigl((g^{-1}w_E)|_L/v\bigr)
  =[I_E:I_E\cap gHg^{-1}]>1.
\]
For every such \(g\), the center of \((g^{-1}w_E)|_L\) on
\(\overline X\) lies in \(\overline X\setminus X\).
\end{proposition}

\begin{proof}
Because \(N\) is the normal closure of \(L/K\), the action of \(G\) on
\(G/H\) is faithful, so \(\bigcap_{g\in G}gHg^{-1}=1\).  Since
\(I_E\neq1\), it is not contained in every conjugate of \(H\); the usual
inertia formula gives the displayed index for a suitable \(g\).  If the
corresponding center belonged to \(X\), the \'etale morphism \(X\to Y\)
would force its ramification index to be one, a contradiction.
\end{proof}

On the generic-degree-three counterexample locus, \(G\cong S_3\); hence
every ramified divisor of the normal closure is visible on some conjugate
sheet, necessarily at its boundary.  Paper~II develops this boundary
geometry for planar Keller maps toward the complementary goal
\(\Obs(F)=0\), equivalently \(R_F=\Delta\).

\section{Conclusion}
\label{sec:conclusion}

The collision construction encodes injectivity by the ideal
\(\Obs(F)\subseteq C_F\).  This ideal is the kernel of diagonal evaluation,
and its vanishing is equivalent to the scheme-theoretic equality
\(R_F=\Delta\).
After passage to the generic fiber, that diagonal becomes the
multiplication factor of \(L\otimes_KL\), while the remaining factors
record the other embeddings of \(L\).

In generic degree three, the diagonal persists as the first factor beside
a quadratic residual algebra, and nonautomorphy is automatic.  Thus the
counterexample condition is equivalent to the Keller condition.  On this
locus the classical Galois case rules out the split, normal branch; hence
the residual algebra is connected, it is the normal closure \(N\), its
Galois group is \(S_3\), and
\[
  K\otimes_BC_F\simeq_K L\times N.
\]
Its nonzero residual factor gives \(\Obs(F)\neq0\).  Hence every map in
the generic-degree-three counterexample locus has the normal-closure factor
\(N\), with \(\Gal(N/K)\cong S_3\).  Counterexamples of other generic degrees
lie outside the present analysis.  The normalization principle
of Section~\ref{sec:further-directions} explains how the same conjugate
sheet geometry leads to boundary questions, which are reserved for the
planar sequel.

\paragraph{Motivating example.}
The explicit three-dimensional polynomial counterexample announced by
Levent Alpöge \cite{Alpoge2026} and independently formalized in
Isabelle/HOL by Ramos, Hulak, and de Queiroz
\cite{AFPJacobianCounterexample} motivated the cubic \(S_3\) comparison.
The formal verification establishes a
constant nonzero Jacobian determinant and distinct rational points with
the same image.  That particular map is not used in the arguments above:
the dimension-three result describes the structure of the
generic-degree-three counterexample locus.

The formalization uses Lean~4 and Mathlib \cite{Lean4,Mathlib}.

\begin{corollary}[First possible generic degree]
\label{cor:first-possible-generic-degree}
If \(F\) is Keller and \([L:K]\leq2\), then \(F\) is a polynomial
automorphism.  Consequently, every complex Jacobian-conjecture counterexample
has generic degree at least three.
\end{corollary}

\begin{proof}
Proposition~\ref{prop:keller-generic-geometry} makes \(L/K\) finite and
separable.  In degree one it is trivial, and in degree two it is Galois.
Theorem~\ref{thm:classical-galois-rigidity} applies in either case.
\end{proof}

This low-degree boundary is the handoff to the planar sequel.  Paper~II asks
for the stronger conclusion that every planar Keller map has generic degree
one---equivalently, by Galois rigidity and the automorphism criterion, that
its collision scheme collapses to the diagonal.

\paragraph{Acknowledgment.}
OpenAI Codex (GPT-5.6 Sol) assisted with the Lean formalization and
manuscript.  The
author reviewed and takes responsibility for all statements, proofs,
citations, and formalization code.

\begingroup
\small
\bibliographystyle{alpha}
\bibliography{references}
\endgroup

\end{document}
